\section{Conclusion}

\paragraph{What worked.} A single BDI pipeline serves both agents. On
Challenge~1 a deliberately simple \texttt{greedy-nearest} strategy is the
strongest or statistically tied on all eight maps, and --- more usefully than
the score itself --- the metrics pipeline let us \emph{diagnose why} and falsify
our own batching hypothesis (Section~\ref{sec:experiments}). On Challenge~2,
Agent~B is genuinely LLM-driven for levels~1--2 (8/8 scenarios green on the
local server, seven by real LLM interpretation) and made latency-safe by the
synchronous pre-apply of safety-critical constraints. Coordination adds claim
deconfliction, mission propagation, and a working one-pickup--another-deliver
handover plus a go-to-and-wait mission. PDDL is integrated as a meaningful but
bounded plan-library member, with a measured, honest account of where it helps
and where the online-solver latency makes BFS the better runtime choice
(Section~\ref{sec:pddl}).

\paragraph{Limitations and future work.} The agent does not model crates or
pushable obstacles, so crate maps (including the live \texttt{crates\_one\_way})
are out of scope. Challenge~2 validation is one run per scenario --- encouraging,
not a statistic, though the failure modes are deterministic in code. The
\texttt{delivery-threshold} strategy under-batches because its early-delivery
penalty is too weak and its threshold is map-independent; auto-tuning the
threshold from the scenario configuration is a natural next step, as is opponent
modeling (dropping contested parcels when an opponent is closer) and a richer
multi-parcel PDDL delivery domain with explicit capacity.

\paragraph{Critical reflection.} The most useful lesson was that the
\emph{infrastructure}, not any single clever heuristic, carried the project:
deterministic low-level movement, a swappable strategy layer, and a metrics
pipeline made it cheap to test ideas and, crucially, to be proven wrong by our
own logs. The strongest Challenge~1 strategy is also the simplest, which is a
result about the \emph{benchmark} (uniform rewards, frequent spawns reward raw
throughput) as much as about the agent; we expect greedy to lose precisely where
Challenge~2 lives --- under missions, constraints, and coordination, where
compliance rather than farming decides the score. The recurring theme across all
three components --- the phantom-carry ack fix, the LLM latency pre-apply, the
PDDL gating --- is the same: in a partially observed, latency-bound, dynamic
world, the safe default is to act deterministically and fast, and to let the
slower, smarter layer (value optimization, the LLM, the symbolic planner) refine
that default only when it is genuinely worth the wait.
